\section{Five Principles of Distributive Justice}
\label{sec:principles}

This section states the five principles in words before any mathematics, because the gap
between the two is the subject of the paper; the wording is the wording participants see.
Each is given as (i) its general claim about what a fair distribution is, (ii) the
assumption that claim needs to be well posed, and (iii) its reading in this domain. Two
features of the domain matter throughout: the thing distributed is a burden rather than a
benefit, so every principle must be read in its harm-allocating form; and the burden is
not freely divisible, since aircraft must fly \emph{somewhere} and the map fixes which
communities lie between the vertiports. A principle here selects among routings, not
among arbitrary allocations.

\subsection{Equal}
\label{sec:p-equal}
\emph{Everyone receives the same share, regardless of differences.} Fairness is
identical treatment: no one gets more or less than anyone else. This assumes all
parties have equivalent standing, and that circumstances are comparable enough for
equal shares to make sense~\cite{dworkin1981}.

\emph{In this domain:} the extra noise added for every region should be as equal as
possible.

Note the form: a \emph{comparative} requirement, constraining differences between
communities rather than any community's absolute level, and phrased as an optimisation
(``as equal as possible'') rather than a threshold. Section~\ref{sec:mismatch} shows
this distinction is what separates Equal from Sufficientarian in principle, and what
makes them collapse into one another in our implementation.

\subsection{Utilitarian}
\label{sec:p-util}
\emph{Maximise overall benefit, or minimise overall harm, across the whole
population.} The best distribution produces the greatest total good even if some
individuals bear more burden than others. This assumes benefits and harms are
measurable and interpersonally comparable, and that aggregate welfare is what
matters~\cite{mill1863,harsanyi1955}.

\emph{In this domain:} the sum of added noise over all regions should be as low as
possible.

Utilitarianism is indifferent to distribution by construction: a routing adding
\SI{1}{\dB} to twenty communities and one adding \SI{20}{\dB} to a single community score
identically.

\subsection{Sufficientarian}
\label{sec:p-suff}
\emph{Ensure everyone meets a minimum threshold before addressing other concerns.} The
priority is that no one falls below an acceptable baseline; above it, further
differences carry no special weight. This assumes there is a meaningful ``enough'' that
can be specified, and that being below the threshold differs in kind from being above
it~\cite{frankfurt1987,crisp2003}.

\emph{In this domain:} no single region's noise level may exceed a stated maximum.
The threshold is a policy input rather than a derived quantity; municipal ordinance
provides one natural anchor, with general residential levels in Austin, TX of around
\SI{70}{\dB} by day and \SI{55}{\dB} at night~\cite{austinnoise}. Note that an ordinance
of this kind bounds the \emph{absolute} level, whereas our implementation bounds the
\emph{added} level --- a substitution whose consequences Section~\ref{sec:concavity}
measures.

Sufficientarianism is the only one of the five natively a \emph{satisfaction} condition
rather than an optimisation: it says when a distribution is acceptable, and is silent
among acceptable ones.

\subsection{Prioritarian}
\label{sec:p-prior}
\emph{Give extra weight to improving the situation of the worst-off.} A benefit matters
more when it goes to someone already disadvantaged, so protecting those who bear the
greatest existing burdens takes priority. This assumes the worst-off can be identified,
and that improving their position is of heightened value~\cite{parfit1997,rawls1971}.

\emph{In this domain:} regions that presently and historically bear the greatest noise
burden from existing transportation infrastructure must receive the least additional
noise.

Prioritarianism requires an input from outside the noise model --- a judgement about who
is already worst-off --- which we make explicitly by naming a protected set of
communities. It is also the principle most exposed to a coupling the source material
flags directly: proximity to a vertiport confers both the worst takeoff and landing
noise \emph{and} the best access to the new service, so removing a burden from a
transit-dependent community may remove a genuine benefit with it.

\subsection{Proportional}
\label{sec:p-prop}
\emph{Those who benefit more should also bear more of the burden.} Burdens and benefits
should be linked rather than distributed independently. This assumes the two can be
meaningfully compared, and that linking them is fairer than treating them
separately~\cite{aristotleNE,young1994,sen1979}.

\emph{In this domain:} areas that gain most from UAM --- in transit access, economic
activity and jobs --- should accept a proportional share of the noise.

Proportional is the only principle whose statement refers to a quantity the noise model
does not contain. ``Benefit'' must be constructed before the principle can be applied,
and every construction embeds contestable judgements about what counts as gaining from an
air-taxi network; Section~\ref{sec:proportional-encoding} makes ours explicit and
participant-editable rather than a hidden constant.

Two independent distinctions among the five govern everything that follows, and it is
worth separating them because they are easily conflated.

The first is \emph{what quantity} a principle speaks about. Sufficientarian and
Prioritarian are \emph{absolute}: they constrain or minimise a community's own noise,
with no reference to anyone else's. Equal is \emph{comparative}: it speaks only about
differences between communities, and is indifferent to the level at which those
differences sit. Utilitarian is \emph{aggregate}, and Proportional is
\emph{relational} --- it ties each community's burden to a separate quantity, its
benefit.

The second is \emph{what form} the principle takes. Utilitarian, Prioritarian and
Proportional rank feasible distributions and prefer the best; they are optimisations.
Sufficientarian partitions distributions into acceptable and unacceptable and is silent
within each class; it is a satisfaction condition. Equal admits \emph{both} readings, and
this is the crux of Section~\ref{sec:mismatch}: ``as equal as possible'' is an
optimisation over the spread, while ``no two communities' increases differ by more than
$X$'' is a satisfaction condition on it. The two are not equivalent, ordinary usage does
not distinguish them, and our implementation --- like most --- takes the second.

The distinction matters because an SMT solver treats the two forms by different
machinery, at different cost and with different behaviour when the clock runs out; and
because, as we show, a comparative satisfaction condition on a non-negative quantity is
weaker than it looks.
